\section{Introduction}
\label{sec:introduction}

Paper C.1 characterized two freight O-D pairs in detail: New York to Los Angeles via the northern transcontinental corridor and Houston to Chicago via the Gulf-Midwest spine \citep{ROUTE_C1}. Both analyses surfaced binding capacity constraints---Donner Pass at 91,200 vehicles per day on a four-lane facility, and the Dallas interchange complex where I-35 and I-45 merge at a sustained V/C of 1.9+---but treated those constraints as local phenomena. A shipper routing from Boston to Portland or from San Antonio to Detroit faces the same bottleneck arcs through a different entry angle. The question C.1 could not answer is whether these two constrained nodes are also the national binding constraints or whether additional, unexamined arcs bind first when flow from all origin clusters is considered simultaneously.

This paper answers that question. We extend the O-D approach of C.1 to national scale by applying the Edmonds-Karp max-flow algorithm to the full ROUTE directed graph: 227 classified corridors, approximately 48,000 directed edges, and eight major freight metro clusters as source-sink terminals \citep{FAF5_2023}. The result is a national bottleneck map grounded in the same network-flow theory that governs single-corridor analysis, but operating at the scale needed to reveal system-level chokepoints that O-D analysis obscures.

The practical stakes are large. The FHWA Freight Analysis Framework projects US truck freight tonnage growing at 1.8 percent per year through 2050, implying a 42 percent increase over the 2023 baseline \citep{FHWA_freight2023}. The current network maximum capacity of 4.8 million commercial vehicles per day leaves no headroom for that growth on the constrained arcs. The gap between current capacity and the 7.2 million vehicles per day required by 2050 demand is not uniformly distributed; it concentrates at precisely the arcs the max-flow minimum-cut theorem identifies as binding. Knowing which arcs bind---and by how much---is the prerequisite for any investment prioritization that claims to be analytically grounded rather than politically driven.

\subsection*{Contributions}

\begin{enumerate}
  \item \textbf{National bottleneck identification.} We apply Edmonds-Karp max-flow to the 227-corridor ROUTE graph and identify the three nationally binding arcs: the I-95 Baltimore-Washington segment (V/C 2.1+), the I-80 Donner Pass crossing (V/C 0.82 at design maximum), and the Dallas I-35/I-45 interchange complex (V/C 1.9+). These three arcs account for the binding minimum cuts in four of six major cross-regional O-D cluster pairs.

  \item \textbf{Incident simulation at national scale.} We simulate two multi-day incidents---a Donner Pass closure and an I-35 Oklahoma corridor disruption---and measure throughput loss and redistribution at the network level. Donner closure drops Northeast-to-Pacific max-flow by 23 percent; the I-40 alternate absorbs the redirected flow and reaches V/C 0.84, within 12,000 vpd of saturation. The simultaneous scenario creates a network failure condition on I-40.

  \item \textbf{Missing link capacity validation.} We add proposed edges for I-69 (Texarkana to Chicago) and I-70 West (Cove Fort, UT to Sacramento, CA) to the ROUTE graph and recompute max-flow. I-69 completion increases Gulf-to-Midwest max-flow by 18 percent by bypassing the Dallas binding constraint. I-70W contributes modestly to baseline throughput (+4.7 percent Northeast-to-Pacific) but reduces Donner-closure throughput loss from 23 percent to 9 percent, making its primary value a resilience asset rather than a capacity asset.
\end{enumerate}

Section~\ref{sec:background} situates this work in the network-flow and freight demand literature. Section~\ref{sec:methods} describes graph construction, demand clustering, and the Edmonds-Karp formulation. Section~\ref{sec:bottleneck} identifies the three binding arcs and maps their geographic distribution. Section~\ref{sec:incident} presents the Donner Pass and I-35 incident simulations. Section~\ref{sec:missing} analyzes I-69 and I-70W as missing link additions. Section~\ref{sec:conclusion} states policy implications and identifies the next extension of this analysis in the D-track climate exposure work.
